\documentclass[11pt,a4paper]{article}
\usepackage[utf8]{inputenc}
\usepackage{amsmath,amssymb,amsthm}
\usepackage{listings}
\usepackage{xcolor}
\usepackage[margin=1in]{geometry}
\usepackage{hyperref}

\newtheorem{theorem}{Theorem}
\newtheorem{lemma}[theorem]{Lemma}
\newtheorem{definition}{Definition}

\lstset{
  language=C++,
  basicstyle=\ttfamily\small,
  keywordstyle=\color{blue}\bfseries,
  commentstyle=\color{green!60!black},
  stringstyle=\color{red},
  numbers=left,
  numberstyle=\tiny\color{gray},
  breaklines=true,
  frame=single,
  tabsize=4,
  showstringspaces=false
}

\title{IOI 2004 -- Empodia}
\author{Solution}
\date{}

\begin{document}
\maketitle

\section{Problem Statement Summary}
Given a permutation $a_0, a_1, \ldots, a_{N-1}$ of $\{0, 1, \ldots, N-1\}$,
find all \textbf{maximal empodia}.

\begin{definition}
An interval $[l, r]$ is a \emph{framed interval} if
$\{a_l, a_{l+1}, \ldots, a_r\}$ is a set of consecutive integers.
An \emph{empodio} is a framed interval $[l, r]$ where $a_l$ is the
minimum and $a_r$ is the maximum of the set (ascending type).
An empodio is \emph{maximal} if it is not properly contained in another
empodio.
\end{definition}

\section{Solution}

\subsection{Key Insight}
Define $c_i = a_i - i$. If $[l, r]$ is a framed interval with
$a_l = \min$ and $a_r = \max$, then $a_r - a_l = r - l$, which means
$c_l = c_r$. Conversely, $c_l = c_r$ is a necessary condition.

\begin{lemma}
An interval $[l, r]$ is an ascending empodio if and only if:
\begin{enumerate}
  \item $c_l = c_r$ (same offset),
  \item $a_l < a_r$,
  \item $a_l = \min(a_l, \ldots, a_r)$ and $a_r = \max(a_l, \ldots, a_r)$.
\end{enumerate}
\end{lemma}

\subsection{Algorithm}
\begin{enumerate}
  \item Group indices by their $c$-value. Within each group (sorted by
        index), check consecutive pairs $(l, r)$ for the empodio
        conditions.
  \item Verify condition (3) using a sparse table for $O(1)$ range
        min/max queries (built in $O(N \log N)$).
  \item Collect all valid empodia, then filter for maximality.
\end{enumerate}

\subsection{Maximality Filter}
Sort candidate intervals by $l$ ascending, breaking ties by $r$
descending. Scan left to right, tracking the maximum $r$ seen so far.
An interval is non-maximal (contained in a previous one) if its $r$ does
not exceed the running maximum. Then reverse-sort the survivors by $r$
descending and filter by $l$ similarly.

\section{C++ Implementation}

\begin{lstlisting}
#include <bits/stdc++.h>
using namespace std;

int main() {
    ios::sync_with_stdio(false);
    cin.tie(nullptr);

    int N;
    cin >> N;
    vector<int> a(N);
    for (int i = 0; i < N; i++) cin >> a[i];

    // c[i] = a[i] - i
    vector<int> c(N);
    for (int i = 0; i < N; i++) c[i] = a[i] - i;

    // Group indices by c-value
    map<int, vector<int>> byC;
    for (int i = 0; i < N; i++) byC[c[i]].push_back(i);

    // Sparse table for range min and range max of a[]
    int LOG = 1;
    while ((1 << LOG) <= N) LOG++;
    vector<vector<int>> rmn(LOG, vector<int>(N));
    vector<vector<int>> rmx(LOG, vector<int>(N));
    for (int i = 0; i < N; i++) rmn[0][i] = rmx[0][i] = a[i];
    for (int k = 1; k < LOG; k++)
        for (int i = 0; i + (1 << k) <= N; i++) {
            rmn[k][i] = min(rmn[k-1][i], rmn[k-1][i + (1 << (k-1))]);
            rmx[k][i] = max(rmx[k-1][i], rmx[k-1][i + (1 << (k-1))]);
        }

    auto qmin = [&](int l, int r) {
        int k = __lg(r - l + 1);
        return min(rmn[k][l], rmn[k][r - (1 << k) + 1]);
    };
    auto qmax = [&](int l, int r) {
        int k = __lg(r - l + 1);
        return max(rmx[k][l], rmx[k][r - (1 << k) + 1]);
    };

    // Find all candidate empodia
    vector<pair<int,int>> cands;
    for (auto& [cv, pos] : byC) {
        for (int idx = 0; idx + 1 < (int)pos.size(); idx++) {
            int l = pos[idx], r = pos[idx + 1];
            if (a[l] < a[r] &&
                qmin(l, r) == a[l] && qmax(l, r) == a[r])
                cands.push_back({l, r});
        }
    }

    // --- Maximality filter ---
    // Pass 1: sort by l asc, r desc. Keep intervals with new max r.
    sort(cands.begin(), cands.end(), [](auto& a, auto& b) {
        return a.first != b.first ? a.first < b.first : a.second > b.second;
    });
    cands.erase(unique(cands.begin(), cands.end()), cands.end());

    vector<pair<int,int>> pass1;
    int maxR = -1;
    for (auto& [l, r] : cands) {
        if (r > maxR) {
            pass1.push_back({l, r});
            maxR = r;
        }
    }

    // Pass 2: sort by r desc, l asc. Keep intervals with new min l.
    sort(pass1.begin(), pass1.end(), [](auto& a, auto& b) {
        return a.second != b.second ? a.second > b.second : a.first < b.first;
    });
    vector<pair<int,int>> result;
    int minL = INT_MAX;
    for (auto& [l, r] : pass1) {
        if (l < minL) {
            result.push_back({l, r});
            minL = l;
        }
    }

    sort(result.begin(), result.end());

    cout << result.size() << "\n";
    for (auto& [l, r] : result)
        cout << l << " " << r << "\n";

    return 0;
}
\end{lstlisting}

\section{Complexity Analysis}
\begin{itemize}
  \item \textbf{Sparse table build:} $O(N \log N)$.
  \item \textbf{Candidate enumeration:} $O(N)$ total across all
        $c$-groups (consecutive pairs sum to $N - 1$), each checked in
        $O(1)$.
  \item \textbf{Maximality filter:} $O(C \log C)$ where $C$ is the
        number of candidates.
  \item \textbf{Overall:} $O(N \log N)$.
  \item \textbf{Space:} $O(N \log N)$ for the sparse table.
\end{itemize}

\end{document}
